\setcounter{chapter}{12}
\chapter{Практический инструментарий: как писать и проверять персонажа}

\section{Character sheet: минимальная анкета персонажа}

Минимальная анкета персонажа — не бюрократическая формальность, а компрессия всего аппарата этой книги в проверяемый список. Она должна отвечать на вопросы, установленные по ходу глав 2--12, а не изобретать новые категории.

\begin{itemize}
\item Архетипический вектор (10.1--10.2): значения по affect, role, intimacy, agency, conflict, visual code, voice.
\item Первичный / вторичный / контрастный слой (10.3).
\item Dere-барьер и его формула (глава 3, 12.2--12.3).
\item Девять осей (глава 5) с кратким значением по каждой.
\item Narrative function, основная и вторичная (2.7).
\item Три failure modes, к которым персонаж ближе всего структурно (глава 9) — не потому что персонаж их демонстрирует, а потому что это его личные "грабли", о которых нужно помнить при письме.
\end{itemize}

Character sheet Рин, заполненный по этой анкете, — почти буквально пересказ глав 3, 5, 7 и 12 в табличной форме; смысл анкеты не в новом содержании, а в том, чтобы иметь его в одном месте для быстрой сверки при написании конкретной сцены.
